% =====================================================================
%  RESUMEN Y ABSTRACT (iguales a los de la versión académica)
% =====================================================================
\chapter*{Resumen}
\addcontentsline{toc}{chapter}{Resumen}
% Guía APA del repositorio (p. 13): un solo párrafo sin sangría, de 150 a 250 palabras;
% «Palabras clave:» en cursiva, en un párrafo nuevo con sangría y separadas por comas.
\noindent Esta investigación diseñó un Sistema de Gestión de la Calidad (SGC) basado en la norma ISO~9001:2026, centrado en el control de materiales, para reducir el costo de ejecución de las losas alivianadas con viguetas pretensadas del Edificio de Posgrado de la Universidad Autónoma Tomás Frías, en Potosí. Como el contrato se resolvió tras la primera losa, el estudio, aplicado, cuantitativo y explicativo, siguió un diseño no experimental: reconstruyó la línea base documental de esa losa (1.109,84~\mdos{}) y evaluó el SGC \emph{ex ante} con un modelo de costos y una simulación de Monte Carlo de \MCiter{} iteraciones sobre los 2.234,60~\mdos{} restantes. El diagnóstico halló un cumplimiento del 19,8~\% de los requisitos de la norma, 22 no conformidades, ninguna probeta de la losa, un exceso de plastoformo del 5,01~\% y acero pagado un 16,75~\% por encima del teórico, inconsistencia que debe conciliarse. El SGC comprende un plan de la calidad (ISO~10005), un plan de puntos de inspección, procedimientos de recepción, control y conciliación de materiales, la gestión de no conformidades y el registro de los costos de la calidad. La simulación estimó una reducción media del costo unitario del \rMedia{}, una probabilidad de \PRpos{} de que el costo disminuya y una reducción del costo de no calidad del \CNCreduccionPct{}; la probabilidad de alcanzar el 5~\% fue de \PHG{}. Se aceptaron las hipótesis específicas (menor desperdicio y menor costo de no calidad), pero no la general. Se diseñó, además, el instrumento para validar el SGC por juicio de expertos.

\emph{Palabras clave:} gestión de la calidad, ISO~9001, control de materiales, desperdicio de materiales, costo de no calidad, simulación de Monte Carlo, losa alivianada, viguetas pretensadas

% ---------------------------------------------------------------------
\chapter*{Abstract}
\addcontentsline{toc}{chapter}{Abstract}
\begin{otherlanguage}{english}
\noindent This study designed an ISO~9001:2026-based Quality Management System (QMS), focused on material control, to reduce the construction cost of the lightweight slabs with prestressed joists of the Graduate School Building of Universidad Autónoma Tomás Frías in Potosí, Bolivia. Because the contract was terminated after the first slab had been built, the applied, quantitative and explanatory study followed a non-experimental design: it reconstructed the documentary baseline of that slab (1,109.84~\mdos{}) and evaluated the QMS \emph{ex ante} with a cost model and a Monte Carlo simulation of \MCiterEN{} iterations over the remaining 2,234.60~\mdos{}. The diagnosis found 19.8\% compliance with the standard's requirements, 22 nonconformities, no concrete specimens taken from the slab, 5.01\% more expanded polystyrene blocks than the theoretical quantity, and slab reinforcement paid 16.75\% above the theoretical quantity, an inconsistency that must be reconciled. The QMS comprises an ISO~10005 quality plan, an inspection and test plan, procedures for material receipt, control and quantity reconciliation, nonconformity management, and quality cost recording. The simulation estimated a mean unit cost reduction of \rMediaEN{}, a probability of \PRposEN{} that the cost decreases, and a \CNCreduccionPctEN{} reduction in the cost of poor quality; the probability of reaching the 5\% target was \PHGEN{}. The specific hypotheses (less material waste and lower cost of poor quality) were accepted, but the general hypothesis was not. An instrument for expert validation of the QMS was also designed.

\emph{Keywords:} quality management, ISO~9001, material control, construction waste, cost of poor quality, Monte Carlo simulation, lightweight slab, prestressed joists
\end{otherlanguage}
